%------------------------------------------------------------------------------%
\documentclass[fleqn,utf8,aspectratio=169,12pt]{beamer}

\usepackage{gaal}

\title{Produto Escalar}

%------------------------------------------------------------------------------%
\begin{document}

\addtitle

%------------------------------------------------------------------------------%
\section{Produto escalar}

%------------------------------------------------------------------------------%
\begin{frame}{Produto escalar}
  Dados os vetores $\vec{u}$ por $\vec{v}$ em $\R^n$
  \[
    \vec{u} = (u_1, u_2, \ldots, u_n)
  \]
  \[
    \vec{v} = (v_1, v_2, \ldots, v_n)
  \]

  \pause
  O \emph{produto escalar}
  \pause
  ou \emph{produto interno}
  \pause
  de $\vec{u}$ por $\vec{v}$ é
  \[
    \pause
    \vec{u}\cdot\vec{v}
    \pause
    \;=\; u_1v_1 + u_2v_2 + \cdots + u_nv_n
    \pause
    \;=\; \sum_{i=1}^n u_i v_i
  \]

  \pause
  O resultado do produto escalar é um número real
  \[
    \vec{u}\cdot\vec{v} \in \R
  \]
\end{frame}

%------------------------------------------------------------------------------%
%------------------------------------------------------------------------------%
\begin{question}
  Calcule o produto escalar dos vetores
  \[
    \vec{u} = ( 2, 3 )
    \qquad\qquad
    \vec{v} = ( 4, -1 )
  \]
\end{question}

%------------------------------------------------------------------------------%
\begin{resolution}{Solução}
  \begin{align*}
    \onslide<+->{\vec{u} \cdot \vec{v}}
    \onslide<+->{& = ( 2, 3 ) \cdot ( 4, -1 ) \\}
    \onslide<+->{& = 2 \times 4 + 3(-1) \\}
    \onslide<+->{& = 8 - 3 \\}
    \onslide<+->{& = 5}
  \end{align*}
\end{resolution}

%------------------------------------------------------------------------------%
\endinput

%------------------------------------------------------------------------------%
\begin{question}
  Calcule o produto escalar dos vetores
  \[
    \vec{w} = \vetortri{1}{-2}{3}
    \qquad\qquad
    \vec{v} = \vetortri{4}{0}{-1}
  \]

\end{question}

%------------------------------------------------------------------------------%
\begin{resolution}{Solução}
  \begin{align*}
    \onslide<+->{\vec{w} \cdot \vec{v}}
    \onslide<+->{& = \vetortri{1}{-2}{3} \cdot \vetortri{4}{0}{-1} \\}
    \onslide<+->{& = 1 \times 4 + (-2) 0 + 3(-1) \\}
    \onslide<+->{& = 4 + 0 - 3 \\}
    \onslide<+->{& = 1}
  \end{align*}
\end{resolution}

%------------------------------------------------------------------------------%
\endinput


%------------------------------------------------------------------------------%
\begin{frame}{Calculando a Norma pelo Produto Escalar}
  \onslide<+->{Dado um vetor
  \(
    \vec{u} = ( u_1, \ldots, u_n ) \in \R^n
  \)}

  \onslide<+->{Calculamos}
  \begin{align*}
    \onslide<+->{\vec{u}\cdot\vec{u} = u_1^2 + \cdots + u_n^2     \\}
    \onslide<+->{\norm{\vec{u}}      = \sqrt{u_1^2 + \cdots + u_n^2}}
  \end{align*}

  \onslide<+->{Comparando temos}
  \[
    \onslide<+->{\vec{u}\cdot\vec{u} = \norm{\vec{u}}^2}
    \qquad\qquad\qquad
    \onslide<+->{\norm{\vec{u}} = \sqrt{\vec{u}\cdot\vec{u}}}
  \]
\end{frame}

%------------------------------------------------------------------------------%
\begin{frame}{Morma como Consequência do Produto Escalar}
  A relação
  \[
    \norm{\vec{u}} = \sqrt{\vec{u}\cdot\vec{u}\,}
  \]
  mostra que o produto escalar \emph{determina} a norma
\end{frame}

%------------------------------------------------------------------------------%
%------------------------------------------------------------------------------%
\begin{question}
  Calcule sua norma do vetor utilizando o produto escalar
  \[
    \vec{u} = (2, -3, 6, -1)
  \]
\end{question}

%------------------------------------------------------------------------------%
\begin{resolution}{Solução}
\begin{columns}
\column{0.6\linewidth}
  \begin{align*}
    \onslide<+->{\vec{u} \cdot \vec{u}}
    \onslide<+->{& = (2, -3, 6, -1) \cdot (2, -3, 6, -1) \\}
    \onslide<+->{& = 2^2 + (-3)^2 + 6^2 + (-1)^2 \\}
    \onslide<+->{& = 4 + 9 + 36 + 1 \\}
    \onslide<+->{& = 50}
  \end{align*}
\column{0.4\linewidth}
  \begin{align*}
    \onslide<+->{\norm{\vec{u}}}
    \onslide<+->{& = \sqrt{\vec{u}\cdot\vec{u}} \\}
    \onslide<+->{& = \sqrt{50} \\}
    \onslide<+->{& = 5\sqrt{2}}
  \end{align*}
\end{columns}
\end{resolution}

%------------------------------------------------------------------------------%
\endinput

%------------------------------------------------------------------------------%
\begin{frame}{Propriedades do produto escalar}
  Para vetores $\vec{u}$, $\vec{v}$ e $\vec{w}$ e escalare $\lambda$

  \pause
  \[
    \vec{u}\cdot\vec{v} = \vec{v}\cdot\vec{u}
  \]

  \pause
  \[
    (\vec{u}+\vec{v})\cdot\vec{w}
    \;=\; \vec{u}\cdot\vec{w} + \vec{v}\cdot\vec{w}
  \]

  \pause
  \[
    (\lambda\vec{u})\cdot\vec{v}
    \;=\; \lambda(\vec{u}\cdot\vec{v})
  \]

  \pause
  \[
    \vec{u}\cdot\vec{u}
    \;=\; \norm{\vec{u}}^2
  \]
\end{frame}

%------------------------------------------------------------------------------%
%\input{file}

%------------------------------------------------------------------------------%
\section{Fórmula Alternativa}

%------------------------------------------------------------------------------%
\begin{frame}{Fórmula Alternativa}
  Dados vetores $\vec{u}$ e $\vec{v}$ com a mesma origem

  \pause
  \medskip
  Seja $\theta$ o ângulo entre eles

  \pause
  \medskip
  Queremos relacionar as componentes dos vetores com o ângulo $\theta$
\end{frame}

%------------------------------------------------------------------------------%
\framedgraphic*{Formula Alternativa}{E-05-Produto_escalar-cosseno-1}{}
\framedgraphic*{Formula Alternativa}{E-05-Produto_escalar-cosseno-2}{}
\framedgraphic*{Formula Alternativa}{E-05-Produto_escalar-cosseno-3}{}

%------------------------------------------------------------------------------%
\begin{frame}{Lei dos cossenos}
  Lei dos cossenos
  \[
    c^2 = a^2 + b^2 - 2ab\cos(\theta)
  \]
  \pause
  No triângulo
  \[
    \pause
    a = \norm{\vec{u}}
    \pause\qquad\qquad
    b = \norm{\vec{v}}
    \pause\qquad\qquad
    c = \norm{\vec{u}-\vec{v}}
  \]
  \pause
  Portanto
  \[
    \pause
    \norm{\vec{u}-\vec{v}}^2
    \;=\; \norm{\vec{u}}^2 + \norm{\vec{v}}^2 - 2\norm{\vec{u}}\norm{\vec{v}}\cos(\theta)
  \]
\end{frame}

%------------------------------------------------------------------------------%
\begin{frame}{Fórmula Alternativa}
  \begin{align*}
    \onslide<+->{\norm{\vec{u}-\vec{v}}^2
    & = \norm{\vec{u}}^2 + \norm{\vec{v}}^2 - 2\norm{\vec{u}}\norm{\vec{v}}\cos(\theta) \\}
    \onslide<+->{2\norm{\vec{u}}\norm{\vec{v}}\cos(\theta)
    & = \norm{\vec{u}}^2 + \norm{\vec{v}}^2 - \norm{\vec{u}-\vec{v}}^2 \\}
    \onslide<+->{& = \vec{u}\cdot\vec{u}
      + \vec{v}\cdot\vec{v}
      - (\vec{u}-\vec{v})\cdot(\vec{u}-\vec{v}) \\}
    \onslide<+->{& = \vec{u}\cdot\vec{u} + \vec{v}\cdot\vec{v}
      - \left(
          \vec{u}\cdot\vec{u} - 2\vec{u}\cdot\vec{v} + \vec{v}\cdot\vec{v}
        \right) \\}
    \onslide<+->{& = \vec{u}\cdot\vec{u}
      + \vec{v}\cdot\vec{v}
      - \vec{u}\cdot\vec{u}
      +2\vec{u}\cdot\vec{v}
      - \vec{v}\cdot\vec{v} \\}
    \onslide<+->{& = 2 \vec{u}\cdot\vec{v}  \\}
    \onslide<+->{\norm{\vec{u}}\norm{\vec{v}}\cos(\theta)
    & = \vec{u}\cdot\vec{v}}
  \end{align*}
\end{frame}

%------------------------------------------------------------------------------%
\begin{frame}{Significado Geométrico}
  A fórmula
  \[
    \vec{u}\cdot\vec{v} \;=\; \norm{\vec{u}}\norm{\vec{v}}\cos(\theta)
  \]
  \pause
  mostra que o produto escalar depende
  \pause
  \begin{itemize}
      \item do comprimento de $\vec{u}$ \pause
      \item do comprimento de $\vec{v}$ \pause
      \item do ângulo entre os vetores
  \end{itemize}
\end{frame}

%------------------------------------------------------------------------------%
\begin{frame}{Significado Geométrico}
  Como
  \(
    -1 \leq \cos(\theta) \leq 1
  \)
  temos

  \pause
  \[
    -\norm{\vec{u}}\norm{\vec{v}}
    \;\leq\;
    \vec{u}\cdot\vec{v}
    \;\leq\;
    \norm{\vec{u}}\norm{\vec{v}}
  \]
\end{frame}

%------------------------------------------------------------------------------%
%\input{file}

%------------------------------------------------------------------------------%
%\input{ex}
%\input{ex}
%\input{ex}
%\input{ex}
%\input{ex}
%\input{ex}
%\input{ex}

%------------------------------------------------------------------------------%
%\section{Lista Mínima}
%
%%------------------------------------------------------------------------------%
%\begin{frame}{Lista Mínima}
%  \begin{enumerate}
%    \item
%  \end{enumerate}
%
%  \vfill
%  Atenção: A prova é baseada no livro, não nas apresentações
%\end{frame}

%------------------------------------------------------------------------------%
\end{document}
%------------------------------------------------------------------------------%
